\documentclass[11pt]{article}
\usepackage[a4paper,margin=25mm]{geometry}
\usepackage[T1]{fontenc}
\usepackage{lmodern,amsmath,amssymb,amsthm,booktabs,array,longtable}
\usepackage[colorlinks=true,linkcolor=blue,citecolor=blue,urlcolor=blue]{hyperref}
\usepackage{orcidlink}
\usepackage{microtype}
\newcommand{\doctype}{Preprint}
\newcommand{\docversion}{v0.01}
\newcommand{\docdate}{3 October 2026}
\newcommand{\versiondoi}{10.5281/zenodo.23128663}
\newcommand{\conceptdoi}{10.5281/zenodo.23128662}
\newcommand{\docrepo}{github.com/fsantibanezleal/CAOS\_RESEARCH}
\newcommand{\tr}{\operatorname{tr}}
\newcommand{\Nd}{N_{\mathrm d}}
\newcommand{\R}{\mathbb R}
\newcommand{\C}{\mathbb C}
\newtheorem{theorem}{Theorem}
\newtheorem{lemma}[theorem]{Lemma}
\newtheorem{proposition}[theorem]{Proposition}
\theoremstyle{remark}\newtheorem{remark}[theorem]{Remark}
\newcommand{\headerblock}{\begin{center}\fbox{\parbox{0.94\textwidth}{\small\raggedright
\textbf{\doctype}\quad$\cdot$\quad\docversion, \docdate\quad$\cdot$\quad Licence CC BY 4.0\\[2pt]
CAOS Research open-problems programme\quad$\cdot$\quad Riemann zeta function\\[2pt]
Version DOI: \href{https://doi.org/\versiondoi}{\versiondoi}\\
Concept DOI (always latest): \href{https://doi.org/\conceptdoi}{\conceptdoi}\\[2pt]
Not peer reviewed. Source code, data and computational results:
\href{https://\docrepo}{\texttt{\docrepo}}
}}\end{center}\vspace{4pt}}
\title{Stronger local overlap certificates and distinct zeros of the Riemann zeta function}
\author{Felipe Santibanez-Leal\,\orcidlink{0000-0002-0150-3246}\\
\small CAOS Research program}
\date{}
\begin{document}
\maketitle\headerblock
\begin{abstract}
We derive a lower asymptotic proportion
$\frac{30945470743359}{36955122080000}=0.8373797460706156\ldots$ of distinct nontrivial
zeta zeros, relative to zeros counted with multiplicity, by applying
Knausg\aa rd's mixed-multiplicity Gram argument to Lavery's seven-point
local theorem with unequal gap pressures. The external local theorem is an
explicit dependency; its reported Lean and nanoda checks were not rebuilt
locally. A separate complete eight-gap interval computation proves a stronger
nine-point local inequality at pressure $1/2500$ and gives the independently
replayed bound $\frac{3997934614153}{4775507750000}=0.8371747724947154\ldots$.
The block transfer uses an elementary spectral dichotomy, without a sharp
trace envelope. Independent rational window enclosures, complete closed-cell
input comparisons and extracted-source checks support the computations.
The analytic input is the corrected integrated unconditional pair-correlation
theorem of Baluyot, Goldston, Suriajaya and Turnage-Butterbaugh.
The results are asymptotic; they do not prove the Riemann hypothesis or give
an effective height.
\end{abstract}

\section{Counts, context and results}
Let $N(T)$ count nontrivial zeros $\rho$ of $\zeta(s)$ with
$0<\operatorname{Im}\rho\le T$, including multiplicity, and let $\Nd(T)$
count the distinct points in the same region. Distinctness refers to the whole
critical strip. A multiple zero contributes one to $\Nd$; two zeros symmetric
about the critical line contribute two. A bound for $\Nd/N$ therefore differs
from a bound for simple zeros on the critical line.

The unconditional pair-correlation method supplied the recent simple-zero
advances of Alp\"oge and Furman and Lamzouri \cite{af,lamzouri}.
Knausg\aa rd \cite{knausgard} retained the Gram defect for critical-line zeros
of multiplicity one or two and obtained
\begin{equation}\label{eq:prior}
 \liminf_{T\to\infty}\frac{\Nd(T)}{N(T)}
 \ge \frac{16260119298029}{19426831050000}
 =0.83699288145242\ldots .
\end{equation}
His matrix inequality, threshold lemma, elementary block dichotomy and
counting assembly are the underlying method here. We give the transfer proof
with a vector of gap pressures and apply two stronger local inputs.

\begin{theorem}[Derived consequence, D]\label{thm:main}
Using Lavery's universal seven-point local theorem \cite{lavery} and the
corrected integrated BGSTB pair-correlation theorem \cite{bgstb,correction},
\begin{equation}\label{eq:main}
 \liminf_{T\to\infty}\frac{\Nd(T)}{N(T)}
 \ge q_7:=\frac{30945470743359}{36955122080000}
 =0.8373797460706156\ldots .
\end{equation}
A separate machine-verified local certificate, with the same analytic and
mixed-Gram inputs, gives
\begin{equation}\label{eq:nine}
 \liminf_{T\to\infty}\frac{\Nd(T)}{N(T)}
 \ge q_9:=\frac{3997934614153}{4775507750000}
 =0.8371747724947154\ldots .
\end{equation}
\end{theorem}

The gain in \eqref{eq:main} over \eqref{eq:prior} is approximately
$0.000386865$, or $0.0386865$ percentage points. The nine-point gain is
approximately $0.000181891$. These are numerical improvements within an
attributed method. Quadratic convex lower bounds, pressure summation and the
scalar window identity below are classical. The source comparison is limited
to the primary material available on 3 October 2026 and does not establish
worldwide priority. The nine-point weights and window come from the source
candidate of tawanerguo-cn \cite{tawan}; the stronger target is proved by the
complete computation described in Section~\ref{sec:certificate}.

\section{A transfer for nonnegative vector pressures}
Let $f$ be bounded, even, nonnegative and supported on $[-1/2,1/2]$, with
$\int f=1$. Write
\begin{gather*}
 K(x)=\int_{-1/2}^{1/2} f(t)e^{2\pi ixt}\,dt,\qquad w(x)=K(x)^2
 \quad(x\in\R),\\
 H_f=2-\left(\int f^2+\iint |s-t|f(s)f(t)\,ds\,dt\right).
\end{gather*}
The kernels used below have positive smooth profiles inside this interval.
Assume nonnegative rational weights $a_{ij}$, $0\le i<j\le r$, satisfy
\begin{equation}\label{eq:capacity}
 \sum_{i=0}^{r-s}a_{i,i+s}\le2\qquad(1\le s\le r).
\end{equation}
Suppose nonnegative pressures $b_0,\ldots,b_{r-1}$ and $\delta>0$ give
the universal local inequality
\begin{equation}\label{eq:local}
 \sum_{i<j}a_{ij}w(g_i+\cdots+g_{j-1})+
 \sum_{i=0}^{r-1}b_i g_i\ge\delta
 \qquad(g_0,\ldots,g_{r-1}\ge0).
\end{equation}
Set $B=\sum b_i$. Equal pressures are not required.

\begin{theorem}[Mixed-Gram transfer, D]\label{thm:transfer}
Let \eqref{eq:capacity}--\eqref{eq:local} hold, and let $m>r$ be an integer.
Choose $\tau>0$ and $c$ such that
\begin{align}\label{eq:conditions}
 D:=\delta(m-r)&\le\tau^2,&
 c&\ge\max(1+\tau,2+\tau/2),\\
 a:=D/m&<2,&
 6c-7-c^2-a&\ge0,&4c-2-c^2-2a&\ge0.\nonumber
\end{align}
Under the corrected integrated BGSTB input, with
$\beta=B(m-r)/m$, one has
\begin{equation}\label{eq:transfer}
 \liminf_{T\to\infty}\frac{\Nd(T)}{N(T)}
 \ge\frac{1+H_f-\beta}{2-a}.
\end{equation}
\end{theorem}

\subsection{Block defect and pressure counting}
For $\lambda\ge0$ put
\[
 \Psi_\tau(\lambda)=(\lambda-1)^2-(\lambda-1-\tau)_+^2.
\]
It is nonnegative and convex, and $\Psi_\tau(1)=0$. A Gram matrix $U$
for $m$ ordered points has unit diagonal. Summing \eqref{eq:local} over its
$m-r$ consecutive windows and using \eqref{eq:capacity} gives
\begin{equation}\label{eq:blockenergy}
 \tr(U-I)^2+P_{\mathrm{block}}\ge D,
\end{equation}
where $P_{\mathrm{block}}$ is the sum of the window pressure charges.
If all eigenvalues of $U$ are at most $1+\tau$, then its clipped defect
$\tr\Psi_\tau(U)$ equals $\tr(U-I)^2$. Otherwise one eigenvalue contributes
more than $\tau^2$, while all other terms and pressures are nonnegative.
Thus $D\le\tau^2$ implies
\begin{equation}\label{eq:dichotomy}
 \tr\Psi_\tau(U)+P_{\mathrm{block}}\ge D.
\end{equation}
This is Knausg\aa rd's elementary dichotomy \cite{knausgard}. A sharp
trace-zero envelope is unnecessary at the parameter choices below.

Order $l$ points $y_0<\cdots<y_{l-1}$ and partition them into complete
consecutive $m$-point blocks for each of the $m$ offsets, leaving shorter
endpoint blocks. Convex trace pinching bounds the full Gram defect from
below by the sum of the complete-block defects. Omitted points may be
pinched into singletons, each contributing zero.

Each possible block start occurs once over all offsets, so the total number
of complete blocks is exactly $\max(l-m+1,0)$. A fixed $r$-gap window lies
in at most $m-r$ complete blocks. A fixed gap appears in local windows with
total pressure at most $\sum_i b_i=B$: each pressure position occurs at
most once. Therefore, with span defined as zero for $l\le1$, averaging
\eqref{eq:dichotomy} yields
\begin{equation}\label{eq:globaldefect}
 \tr\Psi_\tau(U_{\mathrm{full}})
 \ge a l-\beta(y_{l-1}-y_0)-D.
\end{equation}
For $l<m$, the right side is nonpositive, so the same assertion follows
from nonnegativity. This argument proves the pressure coefficient for all
nonnegative vectors. Replacing $B$ by $r\min b_i$ would undercount unequal
charges and is not justified.

\subsection{The matrix inequalities and multiplicities}
For $c>0$ define $\phi_c(x)=x^2-(x-c)_+^2$. We recall and prove the
two attributed matrix inequalities needed for the count \cite{knausgard}.

\begin{lemma}[Threshold inequality]\label{lem:threshold}
If $P\succeq0$, $Q$ is Hermitian and has at most $b$ positive eigenvalues,
and $\tr(P+Q)=N$, then
\[
 \tr(P+Q)^2\ge2cN-2c\tr P+\tr\phi_c(P)-c^2b.
\]
\end{lemma}
\begin{proof}
Write $Q=Q_+-Q_-$ with $Q_+Q_-=0$. The cross term
$\tr((P-Q_-)Q_+)=\tr(PQ_+)\ge0$. Completing squares on $Q_+$ gives
$\tr Q_+^2\ge2c\tr Q_+-c^2b$. Substituting
$\tr Q_+=N-\tr P+\tr Q_-$ leaves
\[
 2cN-2c\tr P+\tr P^2+
 \tr Q_-^2-2\tr((P-cI)Q_-)-c^2b.
\]
The last quadratic, minimized over positive semidefinite $Q_-$, is
$-\tr((P-cI)_+^2)$. This proves the assertion.
\end{proof}

\begin{lemma}[Mixed-multiplicity Gram inequality]\label{lem:mixed}
Let $U$ be a Hermitian Gram matrix with unit diagonal, and
$D_0=\operatorname{diag}(d_i)$, $1\le d_i\le2$. If
$c\ge\max(1+\tau,2+\tau/2)$, then
\[
 \tr\phi_c(D_0^{1/2}UD_0^{1/2})-\tr D_0^2
 \ge\tr\Psi_\tau(U).
\]
\end{lemma}
\begin{proof}
For any Hermitian $G,A$ with $A\preceq2cI$, diagonalizing $G$ gives
\[
 \tr(AG)-\tfrac14\tr A^2
 \le\sum_i(A_{ii}\lambda_i-A_{ii}^2/4)
 \le\sum_i\phi_c(\lambda_i)=\tr\phi_c(G).
\]
The scalar inequality follows by completing a square with $A_{ii}\le2c$.
Set $X=U-I$, $C=\min(X,\tau I)$ by spectral calculus,
$M=D_0^{-1/2}CD_0^{-1/2}$ and $A=2D_0+2M$.
Since $C\preceq\tau I$,
$A\preceq2(D_0+\tau D_0^{-1})\preceq2cI$; the convex function
$d+\tau/d$ on $[1,2]$ is bounded by its endpoint maximum.
Substitution with $G=D_0^{1/2}UD_0^{1/2}$ and its unit diagonal yields
\[
 \tr\phi_c(G)-\tr D_0^2
 \ge2\tr(CX)-\tr M^2
 \ge2\tr(CX)-\tr C^2=\tr\Psi_\tau(U).
\]
The middle inequality is entrywise:
$\tr M^2=\sum_{ij}|C_{ij}|^2/(d_i d_j)\le\tr C^2$.
Neither $C$ nor $M$ needs to be positive semidefinite.
\end{proof}

Let $h$ count critical-line points with multiplicity at least three and $k$
count functional-equation pairs off the line. The remaining
$l=\Nd-h-2k$ points have multiplicity one or two. For the smoothed
complex-zero operator constructed below, their contribution $P$ is positive
semidefinite and the remainder $Q$ has at most $h+k$ positive eigenvalues.
The nonzero spectrum of $P$ is that of $D_0^{1/2}UD_0^{1/2}$.
If $n_1,n_2$ count the low points of multiplicity one and two, then
$\tr D_0^2=n_1+4n_2=3\tr P-2l$.

Combining Lemmas~\ref{lem:threshold} and \ref{lem:mixed}, the remaining
mass is at least $3h+2k$. Because $c\ge2$, increasing its multiplicities
only strengthens the lower estimate. One obtains
\begin{equation}\label{eq:counts}
 \tr(P+Q)^2\ge3N-2\Nd+(6c-7-c^2)h+
 (4c-2-c^2)k+\tr\Psi_\tau(U).
\end{equation}
For example a high point of multiplicity $j\ge3$ contributes
$(2c-3)j+2-c^2\ge6c-7-c^2$; an off-line pair of multiplicity
$j\ge1$ contributes $(4c-6)j+4-c^2\ge4c-2-c^2$.

\subsection{The analytic input and order of limits}
Choose even smooth positive cutoff profiles $f_\epsilon=\eta_\epsilon^2$
supported strictly inside $(-1/2,1/2)$, normalized to integral one, and
converging to $f$ in $L^1$ and $L^2$, with uniformly bounded profiles.
Put $L=\log T$ and
$z_\rho=i(\rho-1/2)L/(2\pi)$ for zeros with $0<\operatorname{Im}\rho\le T$.
Let $F_z(u)=\eta_\epsilon(u)e^{2\pi izu}$, using the inner product linear
in its second slot. The operator
\[
 \mathcal A_\epsilon=
 \sum_z\nu_z F_z\langle F_{\bar z},\,\cdot\,\rangle
\]
is Hermitian by the functional equation. Its trace is $N$, and
\[
 E_\epsilon:=\tr\mathcal A_\epsilon^2
 =\sum_{\rho,\rho'}\nu_\rho\nu_{\rho'}
 K_\epsilon(z_\rho-z_{\rho'})^2.
\]
All operators act on a finite-dimensional span. An off-line pair contributes
$\nu(vw^*+wv^*)=\nu((v+w)(v+w)^*-(v-w)(v-w)^*)/2$, which has at most one
positive eigenvalue. Together with the high critical-line points this
justifies the positive-rank bound used in \eqref{eq:counts}.

BGSTB's integrated Lemma~5 \cite{bgstb} sums over complex zeros with weight
$\omega(\rho-\rho')=4/(4-(\rho-\rho')^2)$. Its preservation under the
corrected pointwise error estimate is explicit in Section~3 of
\cite{correction}. Write $R_\epsilon=f_\epsilon*f_\epsilon$. The Fourier
identity
\[
 \widehat{R_\epsilon-R_\epsilon''/(4L^2)}(z)
 =(1+\pi^2z^2/L^2)K_\epsilon(z)^2
\]
cancels that weight at $z=z_\rho-z_{\rho'}$.
Both $R_\epsilon$ and $R_\epsilon''$ are fixed real even integrable
functions, supported strictly inside $[-1,1]$ and Lipschitz at zero.
Applying the lemma separately to them, then multiplying the latter
asymptotic by $1/(4L^2)$, avoids an error estimate for a varying test function.
The result is
\begin{equation}\label{eq:analytic}
 E_\epsilon=(C(f_\epsilon)+o_\epsilon(1))N,
 \qquad C(f_\epsilon)\longrightarrow 2-H_f.
\end{equation}
This analytic theorem is an attributed dependency, not a conclusion of
the finite interval computations.

To use \eqref{eq:globaldefect} for the smoothed matrix, write
$d_\epsilon=\|f_\epsilon-f\|_1$. Every real kernel entry changes by at
most $d_\epsilon$, uniformly in its argument. An $m$-point block changes
by at most $m d_\epsilon$ in operator norm. Since $\Psi_\tau$ is
$2\max(1,\tau)$-Lipschitz on the nonnegative spectrum, its trace changes
by at most $2\max(1,\tau)m^2d_\epsilon$. Summing fixed-size blocks costs
$O_{m,\tau}(d_\epsilon)N$.

Riemann--von Mangoldt bounds the real low-point span by
$TL/(2\pi)=N+o(N)$. Substituting \eqref{eq:globaldefect} into
\eqref{eq:counts} and comparing with \eqref{eq:analytic} gives
\begin{align*}
 (2-a)\Nd\ge{}&(1+H_f-\beta-O_{m,\tau}(d_\epsilon))N\\
 &+(6c-7-c^2-a)h+(4c-2-c^2-2a)k-o_\epsilon(N)-O_m(1).
\end{align*}
Fix the profile, $m,\tau,c$ first, let $T\to\infty$ next, and only then
let $\epsilon\to0$. The nonnegative residuals in \eqref{eq:conditions}
prove Theorem~\ref{thm:transfer}. No growing-block limit or effective height
is used.

\section{Windows and exact parameter choices}\label{sec:windows}
Both profiles have the form
\begin{equation}\label{eq:profile}
 v(s)=\cos(\sqrt2 s)+\sum_{j=1}^J c_j\cos(2\pi js),\qquad
 f(s)=v(s)/I_1,\quad I_1=\sqrt2\sin(1/\sqrt2),
\end{equation}
on $[-1/2,1/2]$, and zero outside. The numerators of $c_j$, all over
$10^9$, are listed in Table~\ref{tab:coeffs}. The nine-point profile is
from \cite{tawan}; the thirteen-term profile is typh's optimized window
used by Lavery \cite{lavery}. Positivity follows from
$v(s)\ge\cos(1/\sqrt2)-\sum |c_j|>0$.

\begin{table}[htbp]\centering\small
\begin{tabular}{r rr}\toprule
$j$ & nine-point coefficient numerator & thirteen-term numerator\\\midrule
1 & 3322500 & 12497829\\
2 & -7609135 & -14424455\\
3 & 1190194 & 2621498\\
4 & -731476 & 7207519\\
5 & -1680572 & -2561956\\
6 & 1141360 & 6199418\\
7 & 0 & -7197280\\
8 & 0 & 4324511\\
9 & 0 & -5325091\\
10 & 0 & -309374\\
11 & 0 & 443802\\
12 & 0 & -310460\\\bottomrule
\end{tabular}
\caption{Cosine coefficients in \eqref{eq:profile}; each entry is divided by
$10^9$. The leading $\cos(\sqrt2 s)$ coefficient is one.}\label{tab:coeffs}
\end{table}

\begin{proposition}[Scalar identity, D]\label{prop:scalar}
With $\theta=1/\sqrt2$ and \eqref{eq:profile},
\begin{equation}\label{eq:hformula}
 H_f=\frac32-\frac{\cos\theta}{\sqrt2\sin\theta}
 -\frac{\sum_{j=1}^Jc_j^2(1/2-1/(4\pi^2j^2))}{2\sin^2\theta}.
\end{equation}
\end{proposition}
\begin{proof}
Let $\mathcal J$ be the integral operator with kernel $|s-t|$ on
$[-1/2,1/2]$. Solving $(\mathcal J\cos(\omega s))''=2\cos(\omega s)$
with its endpoint conditions gives
\[
 \mathcal J\cos(\omega s)=\frac{\sin(\omega/2)}{\omega}
 +\frac{2\cos(\omega/2)}{\omega^2}
 -\frac{2\cos(\omega s)}{\omega^2}.
\]
Consequently $(I+\mathcal J)\cos(\sqrt2 s)$ is constant. Each
$\cos(2\pi js)$ has zero integral, so the cross terms against the
ground profile vanish. Orthogonality removes the other cross terms and gives
$\langle e_j,(I+\mathcal J)e_j\rangle=1/2-1/(4\pi^2j^2)$.
Thus
\[
 \langle v,(I+\mathcal J)v\rangle
 =\sin^2\theta+\sqrt2\sin\theta\cos\theta+
 \sum c_j^2(1/2-1/(4\pi^2j^2)).
\]
Divide by $I_1^2=2\sin^2\theta$ and subtract from two.
\end{proof}

For the thirteen-term profile an exact rational outward enclosure gives
\begin{gather*}
 0.6721710926412995494027803172404056748101\le H_f,\\
 H_f\le0.6721710926412995494027803172404056748102.
\end{gather*}
The calculation brackets $\sqrt2$ with integer square-root inequalities,
$\pi$ with Machin's formula and forty alternating-series terms, and sine
and cosine with twenty-four Taylor terms and explicit remainders. It uses
only rational endpoints. A separate Arb calculation evaluates the scalar
integrals using an ODE-derived formula and native sinc, checking all 169
symmetry comparisons. For the nine-point profile a source-bound Arb
calculation with a full 4096-cell profile comparison gives
$H_f\ge3362285207/5000000000$.

\begin{table}[htbp]\centering\small
\begin{tabular}{lcc}\toprule
 & seven-point attributed input & nine-point completed computation\\\midrule
$r$ & 6 & 8\\
$\delta$ & $39369/5000000$ & $3051/500000$\\
$B$ & $199193/50000000$ & $8/2500$\\
$H$ used & $33608554629/50000000000$ & $3362285207/5000000000$\\
$m$ & 742 & 958\\
$\tau$ & $12043/5000$ & $2409/1000$\\
$c$ & $17043/5000$ & $3409/1000$\\\bottomrule
\end{tabular}
\caption{Exact inputs for Theorem~\ref{thm:transfer}.}\label{tab:parameters}
\end{table}

For the seven-point input, the six pressure numerators over $10^8$ are
\[
(41468,70344,87381,87381,70344,41468).
\]
Lavery's named theorem
\texttt{cert\_AM} proves \eqref{eq:local} on the full nonnegative orthant;
\texttt{Fw\_W7} identifies its functional with the certified expression.
The source completes its half-space certificate by reflection. Its metadata
attributes the optimized window and evaluator to typh and the weighted-point
refinement to Ainta. The reported external Lean/nanoda checks are preserved,
but were not rebuilt locally.

For the seven-point choice,
\[
 a=\frac{905487}{115937500},\quad
 \beta=\frac{4581439}{1159375000},\quad
 \tau^2-D=\frac{155929}{25000000},
\]
and the two multiplicity residuals are
$16929063061/9275000000$ and $2094101/9275000000$, respectively.
The threshold residuals are $0$ and $2043/10000$.
For the nine-point choice $D=57969/10000<5803281/1000000=\tau^2$;
all remaining rational residuals are likewise nonnegative. Direct fraction
arithmetic in \eqref{eq:transfer} proves both constants in
Theorem~\ref{thm:main}.

\section{The complete nine-point local certificate}\label{sec:certificate}
\begin{proposition}[Machine-verified local inequality, MV]\label{prop:nine}
For the nine-point window in Table~\ref{tab:coeffs} and the exact pair
weights in Appendix~\ref{app:weights}, \eqref{eq:local} holds for every
eight-gap vector $g\ge0$, with $b_i=1/2500$ and $\delta=3051/500000$.
\end{proposition}

The computation uses a grid of spacing $1/4000$ and 128-bit Arb
enclosures. Beyond the exact pressure cutoff, the nonnegative pressure
alone proves the inequality. One-body lower bounds remove whole closed
coordinate cells whose pressure plus adjacent overlap already exceeds the
target. The remaining Cartesian domain has 256 initial boxes, distributed
deterministically across 96 shards. Each node either proves a valid lower
bound or splits an integer cell range into disjoint children whose closed
real boxes cover the parent. No unresolved terminal cell is accepted.

Three bounds certify a node: the nonnegative pressure bound, a sum of
closed-cell kernel lower bounds, or a convex lower bound with a positive
Hessian certificate. For the last bound, each pair term contributes a
signed lower coefficient times a rank-one positive matrix. Their sum $L$
satisfies $\nabla^2F(x)\succeq L$ throughout the box. A floating-point
factorization selects candidate boxes only; positive definiteness is proved
by an Arb LDL factorization of the exact dyadic coefficients.

For a midpoint $x_0$ and gradient $g=\nabla F(x_0)$, integrating the
Hessian on each segment gives
\[
 F(x)\ge F(x_0)+g^T(x-x_0)+\tfrac12(x-x_0)^TL(x-x_0)
 \ge F(x_0)-\tfrac12g^TL^{-1}g.
\]
If $L=\mathcal L\mathcal D\mathcal L^T$, solve $\mathcal Lz=g$ and
compute $g^TL^{-1}g=\sum z_i^2/\mathcal D_{ii}$ by enclosed arithmetic.
The unconstrained minimum is a lower bound even if it lies outside the box.
The algorithm also retains the first-order box bound and accepts either
certified bound. Singular or uncertain pivots reject this path.

All 96 shards complete, with 107,752,902 nodes, 53,876,323 splits and
53,876,579 prunes. The initial-box identity
$107752902=2(53876323)+256$ agrees with the binary-tree accounting.
The prune types are 95,807 pressure, 30,730,268 interval and 23,050,504
convex-bound decisions. The independent standard-library checker reconstructs
the domain, shard allocation, source/table bindings and completed tree counts.
Twelve corrupted copies of actual output are rejected, followed by a fresh
check of unchanged originals. Dense rational quadratic controls separately
compare the LDL solve with exact inversion and reject indefinite, singular
and uncertain-pivot inputs.

\subsection{Independent whole-cell table comparison}
Both stored tables, for $w$ and $w''$, have 61,029 cells. Their lower
coefficients are independently checked on every closed cell at 256-bit
precision through native hypergeometric midpoint jets and Fourier Taylor
remainders. This evaluates a different kernel representation but shares
FLINT/Arb; it is not an independent interval-arithmetic library.

The identity $\operatorname{sinc}(z)={}_0F_1(3/2;-z^2/4)$ supplies the
midpoint jets. Write the unnormalized transform as a finite combination of
shifted sinc functions. If $B_0$ is the sum of the absolute profile
coefficients including the leading one, then
$|K_{\rm raw}''|\le\pi^2B_0$. Enclose the raw kernel on a radius-$R$
subcell by its midpoint value, first derivative and this quadratic remainder,
then divide by its positive normalization before squaring. For
$A_0=B_0/I_1$, the global bound
$|w^{(4)}|\le(2\pi)^4A_0^2$ follows by differentiating the Fourier
representation of $K^2$ supported on $[-1,1]$.
Thus a lower enclosure for $w''$ is
\[
 w''(x_0)-|w'''(x_0)|R-\tfrac12(2\pi)^4A_0^2R^2.
\]
Adaptive bisection proves the comparison on all closed subcells, including
endpoints. The full comparison completes in 225.276 seconds, using
158,883 subcell enclosures and maximum bisection depth three. Deliberately
increased table coefficients are rejected.

\section{Reproducibility and limits}
The exact-byte archive for computation EXP-020 contains the completed
tables, reports, checkpoints, source bindings, actual physical source bytes,
the baseline cache and both standard-library checkers. The 233-member ZIP
has 1,854,946 bytes and SHA256
\begin{center}\small\texttt{91be4798774837bc16007a8236a692078e8}\\
\texttt{25ffb1c8dfa476c05d3c733b17879}.\end{center}
CRC, member and byte comparisons pass. After extraction, both the complete
cover checker and exact transfer checker reproduce the same result.
These transport checks do not rerun all interval operations.

The interval execution trusts the frozen Python verifier, python-flint
0.9.0, FLINT/Arb and directed binary64 enclosures. The complete whole-cell
comparison strengthens the input-table evidence; the metadata checker proves
coverage and binding consistency, not each numerical pruning operation.
There is no tangent-free full replay or end-to-end Lean proof for the
nine-point computation. Windows full-worker replay is documented; other
platforms are not claimed to reproduce its path bindings unchanged.

The seven-point consequence in \eqref{eq:main} trusts Lavery's external
local theorem. Its source SHA256 is
\begin{center}\small\texttt{65564079527487fde93b43bb83dd840a}\\
\texttt{768acf9fb8db1f1f015d38c4faf3b7d4}.\end{center}
The Apache-2.0 source and its external verification log and attestation are
preserved. An independent pure-rational implementation rechecks the window,
source functional and exact transfer. It also checks 288 list/block examples
and 183 pair-count examples with the actual unequal pressures; these finite
examples supplement the general proof rather than establishing universality.

The computation identifiers refer to reproducible source and outputs under
\path{problems/number-theory/riemann-hypothesis/experiments/} in
\href{https://github.com/fsantibanezleal/CAOS_RESEARCH}{the source repository}.
EXP-020 supplies the locally replayed certificate; EXP-025 supplies the
vector-pressure implication and independent rational window calculation.
Supporting EXP-022 bounds the output of the fixed nine-point
window/weight transfer family below $0.8374212877970697\ldots$.
That ceiling restricts the method with those fixed inputs, not the true zero
proportion, and does not restrict other windows or weights.

The counting conclusions are asymptotic and give no effective finite height.
They do not prove RH, universal simplicity, a simple-critical-line proportion,
or an improved short-interval onset. Reported external formal checks and
internal recomputation do not establish peer review. Improvements from
incomplete local covers are not used.

\appendix
\section{Exact pair weights}\label{app:weights}
The upper-triangular entries below are numerators; unspecified diagonal and
lower-triangular entries are zero. Each index span has total weight at most
two. For the nine-point table each span mass equals exactly two.

\begin{center}\scriptsize\setlength{\tabcolsep}{3pt}
\begin{tabular}{rrrrrrrrrr}\toprule
$i\backslash j$ & 0 & 1 & 2 & 3 & 4 & 5 & 6 & 7 & 8\\\midrule
0 & -- & 1802576 & 4832031 & 5411933 & 10000000 & 10000000 & 10000000 & 10000000 & 20000000\\
1 & -- & -- & 2694869 & 2295599 & 1780844 & 0 & 0 & 0 & 10000000\\
2 & -- & -- & -- & 2714860 & 0 & 2807223 & 0 & 0 & 10000000\\
3 & -- & -- & -- & -- & 2787695 & 5744740 & 2807223 & 0 & 10000000\\
4 & -- & -- & -- & -- & -- & 2787695 & 0 & 1780844 & 10000000\\
5 & -- & -- & -- & -- & -- & -- & 2714860 & 2295599 & 5411933\\
6 & -- & -- & -- & -- & -- & -- & -- & 2694869 & 4832031\\
7 & -- & -- & -- & -- & -- & -- & -- & -- & 1802576\\
\bottomrule\end{tabular}\\[4pt]
9-point weights: divide each numerator by $10000000$.
\end{center}

\begin{center}\scriptsize\setlength{\tabcolsep}{3pt}
\begin{tabular}{rrrrrrrr}\toprule
$i\backslash j$ & 0 & 1 & 2 & 3 & 4 & 5 & 6\\\midrule
0 & -- & 21269442 & 890898 & 28398464 & 75539688 & 99999999 & 200000000\\
1 & -- & -- & 35896622 & 65972058 & 71601535 & 48920622 & 99999999\\
2 & -- & -- & -- & 42833935 & 66274085 & 71601535 & 75539688\\
3 & -- & -- & -- & -- & 42833935 & 65972058 & 28398464\\
4 & -- & -- & -- & -- & -- & 35896622 & 890898\\
5 & -- & -- & -- & -- & -- & -- & 21269442\\
\bottomrule\end{tabular}\\[4pt]
7-point weights: divide each numerator by $100000000$.
\end{center}

\begin{thebibliography}{9}
\bibitem{bgstb} S. A. C. Baluyot, D. A. Goldston, A. I. Suriajaya and
C. L. Turnage-Butterbaugh, \emph{An unconditional Montgomery theorem for pair
correlation of zeros of the Riemann zeta-function}, Acta Arith. 214 (2024),
357--376, \href{https://arxiv.org/abs/2306.04799v1}{arXiv:2306.04799v1},
especially Lemma~5.
\bibitem{correction} The same authors, \emph{Pair Correlation of Zeros of the
Riemann Zeta Function I: Proportions of Simple Zeros and Critical Zeros},
\href{https://arxiv.org/abs/2501.14545v3}{arXiv:2501.14545v3}, Section~3.
\bibitem{knausgard} K. M. Knausg\aa rd, \emph{More than 83.69\% of the zeros
of the Riemann zeta function are distinct},
\href{https://arxiv.org/abs/2609.33043v1}{arXiv:2609.33043v1} (2026).
\bibitem{lavery} S. Lavery, \emph{Attempt-013: weighted seven-point refinement
on the optimized AM-form window}, source submission, 3 October 2026,
\href{https://www.riemannzeta.fun/submissions/attempt-013}{riemannzeta.fun/attempt-013}.
Pinned repository revision
\href{https://github.com/josusanmartin/riemann/tree/2fcb7dba28691632cdb0846d694c9b9902cfce66/submissions/attempt-013}{2fcb7dba2869}.
\bibitem{tawan} tawanerguo-cn, \emph{A certified unconditional lower bound for
simple zeros of the Riemann zeta function}, source repository,
\href{https://github.com/tawanerguo-cn/zeta-simple-zeros/tree/45149f6d403059a71be73c5e3f884cee7cd62b20}{revision 45149f6d4030} (2026).
\bibitem{af} L. Alp\"oge and R. Furman, \emph{More than two thirds of the zeta
zeros are simple and on the critical line},
\href{https://arxiv.org/abs/2608.13637}{arXiv:2608.13637} (2026).
\bibitem{lamzouri} Y. Lamzouri, \emph{A new proof that more than 2/3 of the
zeros of the Riemann zeta function are simple and on the critical line},
\href{https://arxiv.org/abs/2609.02882}{arXiv:2609.02882} (2026).
\end{thebibliography}
\end{document}
